\section{용어와 사용 도구}
본문에서 \dag\ 표시가 붙은 용어는 이 부록에 설명이 있으며, PDF에서 용어를 누르면 해당 항목으로 이동한다.
\subsection*{A.1 시뮬레이션과 장면 표현}
\begin{table}[H]
\centering\small
\caption{시뮬레이션·장면 표현 관련 도구}
\begin{tabularx}{\linewidth}{@{}>{\raggedright\arraybackslash}p{30mm}Y@{}}
\toprule
용어 & 설명과 본 과제에서의 역할 \\
\midrule
\glt{gp}{장면 파라미터 (scene parameter)} & 테스트 환경을 정의하는 수치 벡터다. 목표물·장애물·목적지·사람 등 객체의 위치·크기·회전·물성으로 구성되며, 각 항목이 하나의 장면 변수가 된다. 실험 1은 10개, 실험 2는 21개 변수를 사용했다. \\
\glt{envgen}{환경 생성기 (environment generator)} & 장면 파라미터를 입력받아 3D 장면 그래프, 이동 지도, 물리 시뮬레이션 모델을 같은 값에서 함께 생성하는 모듈이다. 세 산출물이 서로 일치하는지, 객체가 겹치거나 허용 범위를 벗어나지 않는지 실행 전에 검사하여 유효하지 않은 환경은 실행하지 않는다. \\
\glt{simmodel}{물리 시뮬레이션 모델 (physics simulation model)} & 물리 엔진이 계산에 사용하는 장면 모델로, 로봇·물체·장애물의 형상, 관절, 질량, 마찰 등을 담는다. 본 과제에서는 MuJoCo의 MJCF 파일로 생성했다. \\
\glt{mujoco}{MuJoCo~\cite{mujoco}} & 접촉이 많은 로봇 동작을 빠르고 정확하게 계산하는 물리 엔진이다. 현재 Google DeepMind가 오픈소스로 관리한다. 두 실험 모두 로봇·물체·장애물의 움직임과 접촉을 MuJoCo로 계산했다. \\
\glt{mjcf}{MJCF (MuJoCo XML)~\cite{mjcf}} & MuJoCo가 읽는 XML 형식의 모델 파일이다. 물체(body), 관절(joint), 충돌 형상(geom), 구동기(actuator), 마찰·질량 같은 물성을 기술한다. 환경 생성기는 장면 파라미터가 바뀔 때마다 이 파일을 새로 만들어 장애물의 위치·크기·회전과 물성을 반영한다. \\
\glt{sg}{3D 장면 그래프 (3D scene graph)~\cite{sg}} & 장면의 객체를 노드, 객체 사이의 공간 관계(거리, 인접, 점유 등)를 간선으로 나타낸 그래프다. 본 과제에서는 객체의 위치·크기·역할과 목표 영역 점유 여부를 구조화하여, 생성된 시뮬레이션 모델과 이동 지도가 서로 일치하는지 검사하는 기준으로 사용했다. 또한 객체 간 거리·목표 영역 점유 같은 관계를 텍스트로 변환하여 AFS의 LLM 입력으로 제공했다. \\
\glt{map}{이동 지도 (occupancy grid map)} & 휴머노이드의 경로 계획에 쓰는 2차원 점유 격자 지도다. 바닥을 일정 크기의 칸으로 나누고 장애물과 상자가 차지하는 칸을 이동 불가로 표시하며, 장면 그래프와 같은 파라미터에서 생성된다. 로봇이 따라갈 궤적 자체가 아니라, 경로 계획기가 궤적을 계산할 때 참조하는 공간 정보다. \\
\glt{robosuite}{robosuite~\cite{robosuite}} & MuJoCo 위에서 로봇팔 조작 환경을 구성하는 시뮬레이션 프레임워크다. LIBERO가 이를 기반으로 만들어졌다. \\
\bottomrule
\end{tabularx}
\end{table}

\subsection*{A.2 로봇과 정책}
\begin{table}[H]
\centering\small
\caption{로봇·정책 관련 도구}
\begin{tabularx}{\linewidth}{@{}>{\raggedright\arraybackslash}p{30mm}Y@{}}
\toprule
용어 & 설명과 본 과제에서의 역할 \\
\midrule
\glt{amr}{AMR} & Autonomous Mobile Robot(자율 이동 로봇). LiDAR·카메라 등 센서로 주변을 인식하고 스스로 경로를 계획해 이동하는 로봇으로, 바닥의 선이나 마커를 따라가는 AGV와 구별된다. 물류·제조 현장에서 물품 운반에 주로 쓰이며, 서론에서 Physical AI 도입 사례로 언급했다. \\
\glt{vla}{VLA~\cite{rt2}} & Vision-Language-Action 모델. 카메라 영상과 자연어 명령을 입력받아 로봇의 관절·그리퍼 동작을 직접 출력하는 신경망 정책이다. 실험 1의 로봇팔 정책이 이에 해당한다. \\
\glt{libero}{LIBERO~\cite{libero}} & 언어 명령 기반 로봇팔 조작 벤치마크로, 130개 과제가 Spatial·Object·Goal·Long 등의 묶음으로 구성된다. 본 과제는 같은 물체를 서로 다른 배치에 두고 ``그릇을 접시에 놓기''를 수행하는 LIBERO-Spatial의 과제 하나를 사용했고, 성공 여부는 LIBERO의 기본 판정을 그대로 따랐다. \\
\glt{panda}{Franka Emika Panda~\cite{panda}} & 연구용으로 널리 쓰이는 7자유도 로봇팔이다. LIBERO 환경의 기본 로봇이며, 실험 1에서 평행 그리퍼로 그릇을 집는다. \\
\glt{groot}{GR00T~N1 / N1.7~\cite{groot,grootrepo}} & NVIDIA가 공개한 로봇용 VLA 기반 모델이다. 영상과 언어를 이해하는 시각-언어 모듈과 연속 동작을 생성하는 동작 생성 모듈로 구성된다. 본 과제는 학습이나 수정 없이 N1.7 버전의 LIBERO-Spatial용 공개 체크포인트를 그대로 사용했다. 한 번의 추론으로 8단계 분량의 동작을 출력한다. \\
\glt{g1}{Unitree G1~\cite{g1}} & Unitree사의 키 약 1.3\,m 휴머노이드 로봇이다. 실험 2에서는 29자유도 시뮬레이션 모델을 사용했다. \\
\glt{wbc}{GR00T 보행 정책~\cite{wbc}} & NVIDIA가 공개한 G1용 학습 기반 보행·균형 제어기다. 상위 계층이 정한 이동 속도와 방향을 관절 명령으로 바꿔 넘어지지 않고 걷게 한다. 본 과제에서는 학습 없이 그대로 사용했다. \\
\glt{agent}{LLM Agent (휴머노이드)} & 실험 2의 상위 의사결정기다. 로봇 카메라 영상, 시뮬레이터가 제공하는 정확한 물체·사람·자기 위치, 목표를 받아 경로 계획, 목표점 이동, 직접 이동, 상자 밀기 같은 로봇 기능 중 하나를 골라 호출한다. 한 번의 선택을 ``로봇 판단''이라 하며, 실행당 판단 횟수를 예비 실험은 10회, 본 실험은 20회로 제한했다. \\
\bottomrule
\end{tabularx}
\end{table}

\subsection*{A.3 탐색·평가 용어}
\begin{table}[H]
\centering\small
\caption{탐색·평가 관련 용어}
\begin{tabularx}{\linewidth}{@{}>{\raggedright\arraybackslash}p{30mm}Y@{}}
\toprule
용어 & 설명 \\
\midrule
\glt{luna}{GPT-6 Luna Pro / Luna} & OpenAI의 대형 언어 모델로, OpenRouter API를 통해 호출했다. 실험 1과 휴머노이드 예비 실험은 Luna Pro, 휴머노이드 본 실험(3세트)은 Luna를 AFS와 로봇 LLM Agent에 사용했다. \\
\glt{afs}{AFS} & Active Failure Search. 이전 실행 결과를 근거로 실패 가능성이 높은 다음 테스트 환경을 골라 실행하는 능동 탐색 방법이다. 본 과제에서는 LLM이 가설과 환경 변경값을 제안하고, 별도 검증기가 범위와 남은 시도 횟수를 확인한다. \\
\glt{random}{Random 테스트} & 비교 기준 방법. 같은 허용 범위에서 바꿀 장면 변수와 그 값을 균등하게 무작위로 뽑으며, 이전 결과를 참고하지 않는다. \\
\glt{var}{장면 변수} & 테스트 장면에서 바꿀 수 있는 하나의 수치 항목(예: 장애물 높이, 통로 폭)이다. 실험 1은 10개, 실험 2는 21개다. \\
\glt{fdr}{FDR} & Failure Discovery Rate. 판정 가능한 유효 실행 중 작업에 실패한 실행의 비율이다. \\
\glt{set}{실험 세트 (난수 시드)} & 독립된 실험 반복 단위로, 시뮬레이션과 무작위 추출의 난수 초깃값(시드)으로 구분한다. 같은 세트에서는 두 방법이 같은 조건으로 시작한다. 로봇팔은 시드 661·673·677, 휴머노이드 본 실험은 898499842·1161219702·253785482(결과를 보기 전에 운영체제 난수로 결정), 예비 실험은 74001이다. \\
\glt{step}{제어 단계} & 로봇 정책이 동작 명령을 한 번 내리는 시간 단위다. 실험 1에서는 1단계가 약 0.05초이며, 300단계(약 15초) 안에 과제를 끝내지 못하면 실패다. \\
\glt{stage}{공통 시작 장면 / 탐색 단계} & 각 세트의 처음 실행(로봇팔 2회, 휴머노이드 1회)은 두 방법이 같은 장면을 쓰는 공통 시작 장면이고, 이후 각 방법이 스스로 장면을 고르는 구간이 탐색 단계다. \\
\glt{basis}{근거 실행 / 반증 조건} & AFS가 제안마다 밝히는 항목이다. 근거 실행은 제안의 근거로 삼은 이전 실행이고, 반증 조건은 가설이 틀렸다고 판단할 결과다. \\
\glt{reg}{회귀 테스트} & 정책이나 현장 조건이 바뀐 뒤, 이전에 저장한 성공·실패 사례를 다시 실행해 결과가 유지되는지 확인하는 테스트다. \\
\bottomrule
\end{tabularx}
\end{table}
